% GENERATED FILE. Edit canonical lessons, then run npm run generate:books.
% canonical-source-hash: fnv1a64:b7ee7499fc12a66d
% canonical-lessons: KA-C543-varanda, KA-C543-landra, KA-C543-mette, KA-C543-pithopakarana, KA-C543-oragu

\chapter{A New Home}
\label{ch:ka-a-new-home}
\hlchaptermodality{543}

\begin{chapteropening}
\textbf{By the end of this chapter:} \emph{I can say a veranda, a lantern, a cushion, furniture and to lean back.}

Complete the last lesson of chapter 543: I can say a veranda, a lantern, a cushion, furniture and to lean back.
\end{chapteropening}

\section[\texorpdfstring{\kn{ವರಾಂಡ} (varāṇḍa)}{ವರಾಂಡ (varāṇḍa)}]{\kn{ವರಾಂಡ} (varāṇḍa) — a veranda}

\label{lesson:KA-C543-varanda}



\begin{quote}
Before the new one: say the Kannada for a basement, then the Kannada for a cabinet.
\end{quote}

\subsection*{You'll want to know: \kn{ವರಾಂಡ}}
\textbf{\kn{ವರಾಂಡ}} — \emph{varāṇḍa} — \textquotedblleft{}a veranda\textquotedblright{}.

\begin{quote}
\kn{ಸಂಜೆ} \kn{ಅಜ್ಜ} \kn{ವರಾಂಡದಲ್ಲಿ} \kn{ಕುಳಿತು} \kn{ಪತ್ರಿಕೆ} \kn{ಓದುತ್ತಾರೆ}.
\end{quote}

\begin{quote}
\emph{sañje ajja varāṇḍadalli kuḷitu patrike ōduttāre.}
\end{quote}

\begin{quote}
In the evening grandfather sits on the veranda and reads the paper.
\end{quote}

A roofed space along the front of the house; \kn{ಜಗಲಿ} is the raised one.

\begin{cousinweb}
\textbf{\kn{ವರಾಂಡ}} (\emph{varāṇḍa}) is a loan from English \emph{veranda}.
\end{cousinweb}

\begin{grammarlens}[title={What you've built}]
One more word for describing what you lived and felt.
\end{grammarlens}

\subsection*{Your turn}
\begin{itemize}
\raggedright
  \item \emph{Say it:} \emph{varāṇḍa}
  \item \emph{Say it:} \emph{varāṇḍa}, once more
  \item \emph{Say it:} \emph{sañje ajja varāṇḍadalli kuḷitu patrike ōduttāre.} — In the evening grandfather sits on the veranda and reads the paper.
  \item \emph{Recall:} say the Kannada for a basement, then the Kannada for a cabinet, then say \emph{varāṇḍa} again
\end{itemize}

\begin{tcolorbox}[breakable,colback=teal!4,colframe=teal!35!black,title={Before you move on}]
What does \kn{ವರಾಂಡ} mean? (\textquotedblleft{}A veranda\textquotedblright{}.) Say it once more.
\end{tcolorbox}
\section[\texorpdfstring{\kn{ಲಾಂದ್ರ} (lāndra)}{ಲಾಂದ್ರ (lāndra)}]{\kn{ಲಾಂದ್ರ} (lāndra) — a lantern}

\label{lesson:KA-C543-landra}



\begin{quote}
Before the new one: say the Kannada for a cabinet, then the Kannada for a veranda.
\end{quote}

\subsection*{You'll want to know: \kn{ಲಾಂದ್ರ}}
\textbf{\kn{ಲಾಂದ್ರ}} — \emph{lāndra} — \textquotedblleft{}a lantern\textquotedblright{}.

\begin{quote}
\kn{ಕರೆಂಟ್} \kn{ಹೋದಾಗ} \kn{ಅಜ್ಜಿ} \kn{ಲಾಂದ್ರ} \kn{ಹಚ್ಚಿದರು}.
\end{quote}

\begin{quote}
\emph{kareṇṭ hōdāga ajji lāndra haccidaru.}
\end{quote}

\begin{quote}
When the power went off, grandmother lit the lantern.
\end{quote}

A glass-covered oil lamp with a handle.

\begin{cousinweb}
\textbf{\kn{ಲಾಂದ್ರ}} (\emph{lāndra}) is usually traced to a European word for a lantern, such as English \emph{lantern}; the path into Kannada is not settled.
\end{cousinweb}

\begin{grammarlens}[title={What you've built}]
One more word for describing what you lived and felt.
\end{grammarlens}

\subsection*{Your turn}
\begin{itemize}
\raggedright
  \item \emph{Say it:} \emph{lāndra}
  \item \emph{Say it:} \emph{lāndra}, once more
  \item \emph{Say it:} \emph{kareṇṭ hōdāga ajji lāndra haccidaru.} — When the power went off, grandmother lit the lantern.
  \item \emph{Recall:} say the Kannada for a cabinet, then the Kannada for a veranda, then say \emph{lāndra} again
\end{itemize}

\begin{tcolorbox}[breakable,colback=teal!4,colframe=teal!35!black,title={Before you move on}]
What does \kn{ಲಾಂದ್ರ} mean? (\textquotedblleft{}A lantern\textquotedblright{}.) Say it once more.
\end{tcolorbox}
\section[\texorpdfstring{\kn{ಮೆತ್ತೆ} (mette)}{ಮೆತ್ತೆ (mette)}]{\kn{ಮೆತ್ತೆ} (mette) — a cushion; a soft mattress}

\label{lesson:KA-C543-mette}



\begin{quote}
Before the new one: say the Kannada for a veranda, then the Kannada for a lantern.
\end{quote}

\subsection*{You'll want to know: \kn{ಮೆತ್ತೆ}}
\textbf{\kn{ಮೆತ್ತೆ}} — \emph{mette} — \textquotedblleft{}a cushion; a soft mattress\textquotedblright{}.

\begin{quote}
\kn{ಅಜ್ಜ} \kn{ಸೋಫಾದ} \kn{ಮೇಲೆ} \kn{ಮೆತ್ತೆ} \kn{ಇಟ್ಟು} \kn{ಕುಳಿತರು}.
\end{quote}

\begin{quote}
\emph{ajja sōphāda mēle mette iṭṭu kuḷitaru.}
\end{quote}

\begin{quote}
Grandfather put a cushion on the sofa and sat down.
\end{quote}

Soft like the word \kn{ಮೆತ್ತಗೆ}, softly.

\begin{cousinweb}
\textbf{\kn{ಮೆತ್ತೆ}} (\emph{mette}) is native Dravidian: DEDR 5068 lists Kannada \emph{mette}, \textquotedblleft{}bedding, mattress\textquotedblright{}.
\end{cousinweb}

\begin{grammarlens}[title={What you've built}]
One more word for describing what you lived and felt.
\end{grammarlens}

\subsection*{Your turn}
\begin{itemize}
\raggedright
  \item \emph{Say it:} \emph{mette}
  \item \emph{Say it:} \emph{mette}, once more
  \item \emph{Say it:} \emph{ajja sōphāda mēle mette iṭṭu kuḷitaru.} — Grandfather put a cushion on the sofa and sat down.
  \item \emph{Recall:} say the Kannada for a veranda, then the Kannada for a lantern, then say \emph{mette} again
\end{itemize}

\begin{tcolorbox}[breakable,colback=teal!4,colframe=teal!35!black,title={Before you move on}]
What does \kn{ಮೆತ್ತೆ} mean? (\textquotedblleft{}A cushion; a soft mattress\textquotedblright{}.) Say it once more.
\end{tcolorbox}
\section[\texorpdfstring{\kn{ಪೀಠೋಪಕರಣ} (pīṭhōpakaraṇa)}{ಪೀಠೋಪಕರಣ (pīṭhōpakaraṇa)}]{\kn{ಪೀಠೋಪಕರಣ} (pīṭhōpakaraṇa) — furniture}

\label{lesson:KA-C543-pithopakarana}



\begin{quote}
Before the new one: say the Kannada for a lantern, then the Kannada for a cushion.
\end{quote}

\subsection*{You'll want to know: \kn{ಪೀಠೋಪಕರಣ}}
\textbf{\kn{ಪೀಠೋಪಕರಣ}} — \emph{pīṭhōpakaraṇa} — \textquotedblleft{}furniture\textquotedblright{}.

\begin{quote}
\kn{ಹೊಸ} \kn{ಮನೆಗೆ} \kn{ನಾವು} \kn{ಹೊಸ} \kn{ಪೀಠೋಪಕರಣ} \kn{ತಂದೆವು}.
\end{quote}

\begin{quote}
\emph{hosa manege nāvu hosa pīṭhōpakaraṇa tandevu.}
\end{quote}

\begin{quote}
We brought new furniture for the new house.
\end{quote}

\kn{ಪೀಠ}, a seat + \kn{ಉಪಕರಣ}, an implement.

\begin{cousinweb}
\textbf{\kn{ಪೀಠೋಪಕರಣ}} (\emph{pīṭhōpakaraṇa}) is a modern learned compound built from Sanskrit \emph{pīṭha}, \textquotedblleft{}a stool, a bench\textquotedblright{} (CDIAL 8222), and \emph{upakaraṇa}, \textquotedblleft{}an implement\textquotedblright{} (CDIAL 2138).
\end{cousinweb}

\begin{grammarlens}[title={What you've built}]
One more word for describing what you lived and felt.
\end{grammarlens}

\subsection*{Your turn}
\begin{itemize}
\raggedright
  \item \emph{Say it:} \emph{pīṭhōpakaraṇa}
  \item \emph{Say it:} \emph{pīṭhōpakaraṇa}, once more
  \item \emph{Say it:} \emph{hosa manege nāvu hosa pīṭhōpakaraṇa tandevu.} — We brought new furniture for the new house.
  \item \emph{Recall:} say the Kannada for a lantern, then the Kannada for a cushion, then say \emph{pīṭhōpakaraṇa} again
\end{itemize}

\begin{tcolorbox}[breakable,colback=teal!4,colframe=teal!35!black,title={Before you move on}]
What does \kn{ಪೀಠೋಪಕರಣ} mean? (\textquotedblleft{}Furniture\textquotedblright{}.) Say it once more.
\end{tcolorbox}
\section[\texorpdfstring{\kn{ಒರಗು} (oragu)}{ಒರಗು (oragu)}]{\kn{ಒರಗು} (oragu) — to lean back; to recline}

\label{lesson:KA-C543-oragu}



\begin{quote}
Before the new one: say the Kannada for a cushion, then the Kannada for furniture.
\end{quote}

\subsection*{You'll want to know: \kn{ಒರಗು}}
\textbf{\kn{ಒರಗು}} — \emph{oragu} — \textquotedblleft{}to lean back; to recline\textquotedblright{}.

\begin{quote}
\kn{ಅಜ್ಜ} \kn{ಗೋಡೆಗೆ} \kn{ಒರಗಿ} \kn{ಕುಳಿತರು}.
\end{quote}

\begin{quote}
\emph{ajja gōḍege oragi kuḷitaru.}
\end{quote}

\begin{quote}
Grandfather sat leaning against the wall.
\end{quote}

Past (I …): \kn{ಒರಗಿದೆ}, \emph{oragide}.

Against what: -\kn{ಗೆ}. \kn{ಒರಗು} \kn{ಕುರ್ಚಿ}, an easy chair.

\begin{cousinweb}
\textbf{\kn{ಒರಗು}} (\emph{oragu}) is native Dravidian: DEDR 707 lists Kannada \emph{oṟagu}, \textquotedblleft{}to incline, recline, lean upon, lie down\textquotedblright{}, and the noun \emph{oṟagu}, \textquotedblleft{}a cushion to lean the back against\textquotedblright{}.
\end{cousinweb}

\begin{grammarlens}[title={What you've built}]
One more word for describing what you lived and felt.
\end{grammarlens}

\subsection*{Your turn}
\begin{itemize}
\raggedright
  \item \emph{Say it:} \emph{oragu}
  \item \emph{Say it:} \emph{oragu}, once more
  \item \emph{Say it:} \emph{ajja gōḍege oragi kuḷitaru.} — Grandfather sat leaning against the wall.
  \item \emph{Recall:} say the Kannada for a cushion, then the Kannada for furniture, then say \emph{oragu} again
\end{itemize}

\begin{tcolorbox}[breakable,colback=teal!4,colframe=teal!35!black,title={Before you move on}]
What does \kn{ಒರಗು} mean? (\textquotedblleft{}To lean back; to recline\textquotedblright{}.) Say it once more.
\end{tcolorbox}
